\section{Introduction}
\label{sec:intro}

Large language models are rapidly migrating from short, single-turn dialogue to long-horizon \emph{agentic} workflows: writing and shipping production code, navigating the open web, operating heterogeneous tools, and producing structured office artifacts across hundreds of interleaved reasoning and action steps~\citep{openai2025gpt5,anthropic2025claude46,google2025gemini31}. This shift exposes two distinct difficulties. First, the inherently ultra-long context of agentic tasks introduces formidable efficiency and cost bottlenecks during both training and inference, particularly under the stringent requirements of large-scale, high-availability production deployment. Second, deployment in the wild demands solving intrinsically complex and high-stakes tasks, such as production-grade software engineering and knowledge-intensive office automation.

To address these twin challenges, we introduce the \textbf{MiniMax-M2 series}, a family of Mixture-of-Experts (MoE) language models built around a single design principle: \emph{mini activations can unleash maximum real-world intelligence}. The flagship M2 is a 62-layer decoder-only Transformer with 229.9B total parameters and only 9.8B activated per token, organized as 256 fine-grained experts~\citep{dai2024deepseekmoe} with sigmoid gating, full multi-head attention with GQA~\citep{ainslie2023gqa}, a 192K-token native context window, and a Multi-Token Prediction (MTP) module~\citep{gloeckle2024better,deepseekai2024v3} that doubles as a speculative-decoding draft path~\citep{leviathan2023fast} at inference. Pre-training on 29.2T tokens establishes the base; the bulk of M2's real-world capability is then constructed by an agent-native post-training pipeline whose components co-evolve from M2 through M2.5 to the latest M2.7.

\noindent \textbf{Main contributions.} The continuous capability evolution and performance enhancements of the MiniMax-M2 series stem primarily from the following technical innovations:
\begin{itemize}
    \item We design high-fidelity, large-scale \textbf{agent data pipelines} tailored for agentic coding, collaborative work (cowork), reasoning, and general knowledge tasks, where each task is accompanied by its corresponding static/runtime environments, verifiable rewards, or credible feedback signals. We find that elevating the reward quality and credibility of each accepted trajectory---whether through executable verification signals or judge-model evidence checking---is of paramount importance to fully unleashing the inherent potential of the base model.

    \item We build \textbf{Forge}, an agent-native RL system engineered for large-scale, general-purpose agentic reinforcement learning, which seamlessly admits both white-box and black-box (API-only) agents within a unified training loop. By decoupling key architectural components---including training, inference, and the agent itself---and pairing this separation with robustness-first algorithmic designs and a meticulous reward system, Forge achieves highly stable RL-time scaling. Furthermore, Forge incorporates windowed-FIFO scheduling to absorb trajectory-length variance, prefix-tree merging, and inference kernels co-designed with our deployment stack, thereby substantially boosting RL training efficiency and scalability.

    \item We demonstrate, in M2.7, an early operational form of \textbf{self-evolution}: the model autonomously triages failed training runs on our own infrastructure, edits its own agent scaffold across tasks and experiments, and is evaluated by running multi-round self-improvement on representative ML-engineering tasks. The within-series gains from M2 $\to$ M2.5 $\to$ M2.7 on agentic benchmarks already reflect this, closing one of the most expensive human-in-the-loop bottlenecks in frontier model development.
\end{itemize}

\noindent \textbf{Results.} Figure~\ref{fig:main-bar} previews the headline numbers for MiniMax-M2.7 across three capability areas. On \emph{agentic coding}, M2.7 reaches 56.2 on \emph{SWE-bench Pro}, 76.5 on \emph{SWE-bench Multilingual}, 52.7 on \emph{Multi-SWE-bench}, and 57.0 on \emph{Terminal-Bench 2.0}. On \emph{agentic cowork}, it reaches 62.7 on \emph{MM Claw}, 77.8 on \emph{BrowseComp}, 50.0 on \emph{GDPval-AA}, and 46.3 on \emph{Toolathlon}. On \emph{reasoning \& knowledge}, M2.7 posts 94.2 on \emph{AIME 2026} and 89.8 on \emph{GPQA-Diamond}. With only $\sim$10\,B activated parameters, MiniMax-M2.7 approaches the performance of the strongest closed-weight frontier systems. We refer the reader to \Cref{sec:eval} for the full benchmark suite, per-benchmark analysis, and within-series progression.
